\chapter*{Abstract}
\addcontentsline{toc}{chapter}{Abstract}

This report consolidates every theoretical model, numerical sweep, and validation
result this project has produced on a single question: \emph{for a myopic eye viewing
through a single circular pinhole aperture in an otherwise opaque film, what aperture
diameter is optimal, and by which criterion?} It deliberately excludes the project's
separate multi-hole array, ghosting, packing-geometry, and lens-material work, which
answer a different question (how to \emph{tile} many holes safely) and are documented
elsewhere.

Four largely independent modeling efforts are covered, each with its own governing
equations, numerical implementation, and full-grid sweep across 15 apertures
($75$--$2000\um$) and 8 myopia levels ($0$--$8\,$D): (1) an exact scalar
Angular-Spectrum-Method wave-optics pipeline from pinhole to retina, whose blur-only
optimum is the project's headline number; (2) a brightness-constrained variant of the
same pipeline, which shows that a minimum-illuminance floor moves the optimum by
roughly $2$--$3\times$ into the clinically-observed bracket; (3) a human-vision layer
(MTF, contrast sensitivity, Strehl ratio, Watson--Yellott pupil model, troland
photometry) built on top of the same optics, producing a predicted logMAR acuity
surface; and (4) a structurally independent ideal-observer Monte Carlo acuity model
(R7acuity), built specifically to cross-check (3) with a different architecture.

A significant fraction of this project's actual scientific content is a documented
\emph{physics audit}: a real lens-phase/retina-propagation wavelength bug was found,
independently re-derived two ways, and fixed, moving the blur-only optimum aperture by
up to $+150\%$ at some myopia levels; a separate Python import-shadowing bug meant two
previously-published acuity-validation results had silently never been evaluated
against the corrected optics at all, and both were re-run once the shadowing was
fixed. Both corrections, and their full before/after numbers, are reported here in
full --- including the honest result that neither acuity model achieves statistical
significance against the two available clinical datasets (Wang 2009, $n=5$; Kim et al.
2017, one aggregate point), a negative result this report does not soften.

The synthesis chapters bring these efforts together: an aperture-optimum comparison
across four different blur-percentile definitions ($r_{EE}50\%/70\%/80\%/90\%$),
including a previously-undocumented numerical saturation artifact at low myopia that
is identified and quantified rather than silently averaged over; a full corrected-vs-deprecated
optics comparison; and a cross-model summary table spanning every method in this
report. No single "recommended aperture" is asserted anywhere in this document --- each
chapter reports what its own model predicts, under its own stated assumptions, and the
final synthesis chapter shows how those predictions relate to one another and to the
two independently-documented patents in this space, without collapsing them into one
number.
